\chapter*{Appendix C}
\label{appendix:microfluidic_protocol}

\section{Detailed protocol for microfluidic device fabrication}

\subsection{Scope and device revision}
This appendix describes the fabrication workflow for the two-layer SU-8 master, PDMS casting, glass bonding, and epoxy replica master generation used in this thesis. The protocol is intended to be reproducible by users with prior cleanroom and soft-lithography experience.

\textbf{Device name:} Mother machine for \textit{S. pombe} long-term imaging \\
\textbf{Revision ID used in this thesis:} [fill: e.g., MM-SP-v3.2] \\
\textbf{Date range used:} [fill: e.g., 2024-10 to 2026-01] \\
\textbf{Minimum critical feature size:} [fill] \\
\textbf{Critical dimensions:} observation chamber width and chamber-to-main-channel alignment

\subsection{Design files and critical dimensions}
Store design files with explicit versioning and archive the exact version used in each experiment batch.

\textbf{Required archived files}
\begin{itemize}
  \item CAD layout file (native format)
  \item Mask export (GDSII or writer-native format)
  \item Dimensioned PDF drawing
  \item This protocol text (same revision as mask)
\end{itemize}

\textbf{Geometry used in this work}
\begin{table}[H]
\centering
\caption{Device geometry and tolerance priorities}
\label{tab:device_geometry}
\begin{tabular}{lll}
\hline
\textbf{Feature} & \textbf{Nominal value} & \textbf{Priority} \\
\hline
Main channel width & 180 $\mu$m & Critical \\
Observation chamber width & 6.3 $\mu$m & Critical \\
Observation chamber length & 40 $\mu$m & Critical \\
Port diameter & 0.5 mm punch & Moderate \\
Chip outer outline & [fill] & Non-critical \\
\hline
\end{tabular}
\end{table}

\subsection{Photomask and alignment specification}
\begin{itemize}
  \item Mask fabrication: maskless lithography system (\textmu mLA, Heidelberg Instruments) on mask blanks (Clean Surface Technology).
  \item Tone convention: SU-8 is negative-tone resist; exposed regions define raised master features.
  \item Alignment marks: Layer 2 (drain layer) is aligned to Layer 1 marks.
  \item Alignment acceptance: overlay error $\leq$ [fill] $\mu$m at three representative mark pairs.
\end{itemize}

\subsection{Reagents and equipment}
\begin{table}[H]
\centering
\caption{Core materials and instruments}
\label{tab:fab_materials}
\begin{tabular}{ll}
\hline
\textbf{Category} & \textbf{Item} \\
\hline
Photoresists & SU-8 2, SU-8 3005, SU-8 3010 (MicroChem) \\
Developer & SU-8 developer (MicroChem) \\
Elastomer & SYLGARD 184 (Dow Silicones) \\
Replica material & LOCTITE EA E-30CL (2-part epoxy) \\
Substrate & Silicon wafer (University Wafers) \\
Cover glass & Matsunami NEO No.1 (0.13--0.17 mm) \\
Spin coater & Mikasa MS-A150 \\
Mask aligner & Mikasa MA-20 \\
Plasma system & FA-1 (SAMCO) \\
Punch & 0.5-mm biopsy punch (BP-A05F) \\
\hline
\end{tabular}
\end{table}

\subsection{Master mold fabrication (two-layer SU-8)}

\subsubsection{General notes}
\begin{itemize}
  \item Run process on dehydrated wafers and stable room humidity.
  \item Record resist lot number and room temperature for each batch.
  \item Measure final layer thickness by profilometer at 3--5 positions.
\end{itemize}

\subsubsection{Layer 1: observation channel layer}
\textbf{Action}
\begin{enumerate}
  \item Prepare mixed resist from SU-8 2 and SU-8 3005 at \textbf{explicit ratio} [fill: weight ratio or volume ratio].
  \item Dispense 1 mL resist on silicon wafer.
  \item Spin-coat at 3000 rpm for 30 s (Mikasa MS-A150).
  \item Soft-bake at 65$^{\circ}$C for 1 min, then 95$^{\circ}$C for 2 min.
  \item Align observation-channel mask and expose using MA-20 (3 exposures, 5.0 s each).
  \item Post-exposure bake at 65$^{\circ}$C for 1 min, then 95$^{\circ}$C for 3 min.
  \item Develop in SU-8 developer for [fill] min with [fill: gentle/manual] agitation.
  \item Rinse with isopropanol and dry with nitrogen.
\end{enumerate}

\textbf{Critical notes}
\begin{itemize}
  \item Report exposure as dose (mJ/cm$^2$) if intensity calibration is available.
  \item If only exposure time is available, record aligner model and most recent calibration date.
  \item If edge bead appears, remove before alignment of Layer 2.
\end{itemize}

\textbf{Expected result}
\begin{itemize}
  \item Target thickness: [fill] $\mu$m
  \item Measured thickness: [fill] $\mu$m (mean $\pm$ SD, n = [fill])
  \item No collapsed features or residual resist in observation channels
\end{itemize}

\subsubsection{Layer 2: drain layer}
\textbf{Action}
\begin{enumerate}
  \item Dispense SU-8 3010 on wafer containing completed Layer 1.
  \item Spin-coat at 1500 rpm for 30 s.
  \item Soft-bake at 65$^{\circ}$C for 1 min, then 95$^{\circ}$C for 8 min.
  \item Align drain mask to Layer 1 alignment marks.
  \item Expose for 30 s (record intensity/dose if available).
  \item Post-exposure bake at 65$^{\circ}$C for 3 min, then 95$^{\circ}$C for 10 min.
  \item Develop, rinse, and dry as in Layer 1.
\end{enumerate}

\textbf{Expected result}
\begin{itemize}
  \item Target thickness: [fill] $\mu$m
  \item Overlay error: $\leq$ [fill] $\mu$m
  \item Continuous drain features without residual scum
\end{itemize}

\subsection{PDMS casting, porting, and bonding}

\subsubsection{PDMS casting and demolding}
\begin{enumerate}
  \item Mix SYLGARD 184 base and curing agent at 10:1 ratio (recommended: by weight).
  \item Degas under vacuum for 1 h (or until no visible bubbles $>$ [fill] mm).
  \item Pour over SU-8 master and cure at 65$^{\circ}$C for at least 12 h.
  \item Cut cured PDMS blocks and punch inlet/outlet ports using 0.5-mm punch.
  \item Sonicate PDMS blocks in ethanol and dry at 65$^{\circ}$C.
\end{enumerate}

\subsubsection{Cover glass cleaning}
\begin{enumerate}
  \item Sonicate cover glasses sequentially in diluted Contaminon LS-II, ultrapure water, ethanol, and ultrapure water.
  \item Treat with 0.1 M NaOH.
  \item Rinse with ultrapure water and dry at 140$^{\circ}$C.
\end{enumerate}

\subsubsection{Plasma bonding}
\begin{enumerate}
  \item Treat PDMS and cleaned cover glass with plasma etcher (FA-1, SAMCO) for 10 s.
  \item Bring surfaces into conformal contact immediately after activation.
  \item Apply gentle uniform pressure to avoid trapped air.
\end{enumerate}

\textbf{Critical notes}
\begin{itemize}
  \item Record gas type, RF power, chamber pressure, and flow rate when available.
  \item Define a strict activation-to-contact window (recommended: within 60 s).
\end{itemize}

\textbf{Expected result}
\begin{itemize}
  \item No visible interfacial voids at channel region
  \item Bond survives tubing connection and initial perfusion without delamination
\end{itemize}

\subsection{Epoxy replica master fabrication}
\begin{enumerate}
  \item Cast PDMS negative from silicon master (10:1 SYLGARD 184, degas 1 h, cure 65$^{\circ}$C for $\geq$12 h).
  \item Cast two-part epoxy (LOCTITE EA E-30CL) on flat plastic plate using PDMS negative mold.
  \item Cure for at least 24 h.
  \item Use cured epoxy as reusable master for subsequent PDMS casting.
\end{enumerate}

\textbf{Process controls}
\begin{itemize}
  \item Keep curing surface level to minimize warp.
  \item Record casts per epoxy master and retire master when defects appear.
  \item Store masters in dust-free closed containers.
\end{itemize}

\subsection{Quality control and acceptance criteria}
\begin{table}[H]
\centering
\caption{Batch-level quality control checks}
\label{tab:fab_qc}
\begin{tabular}{lll}
\hline
\textbf{Check} & \textbf{Method} & \textbf{Pass criterion} \\
\hline
Layer thickness & Profilometer (3--5 points) & Within $\pm$[fill]\% of target \\
Feature integrity & Brightfield microscope & No blocked or collapsed channels \\
Alignment & Mark-to-mark overlay & $\leq$ [fill] $\mu$m \\
Bond quality & Visual interface scan & No voids near channels \\
Leak test & Perfusion test & No leakage at 2 mL/h for [fill] min \\
\hline
\end{tabular}
\end{table}

\subsection{Troubleshooting guide}
\begin{table}[H]
\centering
\caption{Common failure modes and corrective actions}
\label{tab:fab_troubleshooting}
\begin{tabular}{p{0.23\linewidth} p{0.33\linewidth} p{0.33\linewidth}}
\hline
\textbf{Symptom} & \textbf{Likely cause} & \textbf{Corrective action} \\
\hline
Channels partially blocked & Underdevelopment or resist residue & Extend develop time; refresh developer; verify agitation \\
Poor Layer 1/2 overlay & Edge bead or mark contrast failure & Remove edge bead; improve mark visibility; reduce handling delay \\
PDMS tears at ports & Dull punch or excessive force & Replace punch; punch on hard backing; clean debris after punching \\
Bonding voids or leaks & Particles, delay after plasma, low activation & Improve cleaning; contact within defined time window; optimize plasma settings \\
Difficulty demolding PDMS & Master surface adhesion too strong & Apply compatible anti-stiction treatment; reduce peeling angle \\
Warped epoxy master & Non-level curing surface or uneven stress & Use leveled base; apply uniform support during cure \\
\hline
\end{tabular}
\end{table}

\subsection{Document control and version history}
Track all process changes that affect dimensions, yield, or bond quality.

\begin{table}[H]
\centering
\caption{Protocol revision history}
\label{tab:fab_revision_history}
\begin{tabular}{llll}
\hline
\textbf{Revision} & \textbf{Date} & \textbf{Change summary} & \textbf{Author} \\
\hline
v1.0 & [fill] & Initial protocol used in thesis & [fill] \\
v1.1 & [fill] & [fill] & [fill] \\
\hline
\end{tabular}
\end{table}

\subsection{Safety note}
All steps involving SU-8 resists, SU-8 developer, NaOH, ethanol, plasma equipment, and epoxy components must follow institutional chemical safety rules and relevant SDS documents. PPE, waste segregation, and local ventilation requirements should be applied according to facility policy.
